The primary output is the interval $[L_\phi,U_\phi]$. We report its width
\[
  W_\phi = U_\phi - L_\phi
\]
as an identified-set diagnostic. Large width means the supplied evidence leaves
substantial composition risk unidentified.

When a selected or observed event risk $r$ is available, we report its position
inside the interval:
\[
  \operatorname{pos}(r) = \frac{r-L_\phi}{U_\phi-L_\phi},
\]
with the degenerate case handled separately. This is descriptive, not a safety
claim.

We also compute the product-coupling baseline $r_{\mathrm{prod}}$ obtained by
treating singleton failures as independent Bernoulli variables. The
independence regret
\[
  r-r_{\mathrm{prod}}
\]
is a signed diagnostic for how far an observed or selected risk lies from that
explicit assumption. The product-coupling baseline is useful as a comparison
point, but it is not the default model of reality.

Every reported sharp interval should include endpoint witness distributions
$\pi_L$ and $\pi_U$. These witnesses satisfy the declared constraints and attain
the lower and upper query values. They make the mathematical result
recomputable and challengeable. They do not prove that the empirical inputs are
representative or that either endpoint is realized in deployment.
